% Generated draft: requires mathematical and source review.
\documentclass[11pt]{article}
\usepackage[margin=25mm]{geometry}
\usepackage{fontspec}
\setmainfont{Latin Modern Roman}
\usepackage{amsmath,amssymb,mathtools,unicode-math}
\setmathfont{Latin Modern Math}
\usepackage{enumitem,needspace,xurl,hyperref}
\setlength{\parindent}{0pt}
\setlength{\parskip}{5pt}
\setlength{\emergencystretch}{3em}
\newcommand{\N}{\mathbb{N}}
\newcommand{\Z}{\mathbb{Z}}
\newcommand{\Q}{\mathbb{Q}}
\newcommand{\R}{\mathbb{R}}
\newcommand{\C}{\mathbb{C}}
\newcommand{\F}{\mathbb{F}}
\newcommand{\A}{\mathbb{A}}
\newcommand{\PP}{\mathbb{P}}
\newcommand{\E}{\mathbb{E}}
\newcommand{\eps}{\varepsilon}
\newcommand{\dd}{\,\mathrm{d}}
\begin{document}
\section*{UnsolvedMath ID: 30000740}
\textbf{Source problem code:} catalogue entry
\subsection{Definitions and mathematical statement}
Let $(V,F)$ be the hypergraph in the source, with $V$ a finite board and $F$ its family of winning sets under the usual finite positional-game interpretation. In the unbiased Maker--Breaker game, players claim previously unclaimed vertices one at a time. Maker wins by claiming every vertex of some member of $F$; Breaker wins by preventing this.

In the standard pairing-game convention, Picker offers two unclaimed vertices, Chooser takes one, and Picker receives the other. In Picker--Chooser, Picker has the Maker objective of completing a member of $F$. In Chooser--Picker, Chooser has that objective and Picker tries to prevent it. The source does not spell out these role conventions or the rule for a last unpaired vertex, so the definitions must be matched with the report rather than inferred from names alone.

The two proposed strategy transfers, for every admissible hypergraph with its fixed winning sets, are
\begin{enumerate}
\item If Maker has a winning Maker--Breaker strategy when Maker plays second, then Picker has a winning Picker--Chooser strategy.
\item If Breaker has a winning Maker--Breaker strategy when Breaker plays second, then Picker has a winning Chooser--Picker strategy.
\end{enumerate}
The second-player condition refers to the named player in the Maker--Breaker premise: the turn order in the two premises is therefore different. The conclusions always concern Picker, whose objective differs between the two pairing games. No converse, first-player premise, or biased version is added.

The original source note reports a proved special case for the Picker--Chooser version of generalized Shannon switching games on matroids. It does not state a theorem for all hypergraphs, all matroid-related games, or both transfer directions. The dated assessment summarizes Shannon-switching/matroid cases more broadly; the exact special-case hypotheses and scope require source review. The general pair of transfer questions is retained separately from that supplied partial result.
\subsection{Short English statement}
The conjectures transfer a second-player Maker win to a Picker win in Picker--Chooser, and a second-player Breaker win to a Picker win in Chooser--Picker. The roles and turn orders are distinct, and the supplied matroid special case is narrower than a general transfer theorem.
\subsection{Sources}
\begin{itemize}
\item Local UnsolvedMath registry, permanent ID 30000740, code catalogue entry, ``Picker--Chooser and Maker--Breaker Strategy Transfers.'' Proposed year 2007; proposer unspecified.
\item \emph{Mini-Workshop: Positional Games} (2008), Oberwolfach Reports, pp.\ 1073--1114, DOI 10.4171/owr/2007/20; \url{https://doi.org/10.4171/owr/2007/20}. The displayed report year 2008 differs from the DOI and proposed year 2007. This is the sole supplied literature reference.
\item The retained original reports the Picker--Chooser generalized-Shannon-switching-on-matroids special case; no independent theorem citation is supplied.
\item Statement status corrected\_verified; the supplied verification says official-report content was checked and wording/notation normalized without changing the source claim.
\item Literature triage checked 2026-08-21: partially solved. It says the general transfer conjecture is not established and the report proves Shannon switching/matroid cases. This summary must be checked against the original's specific pairing-game clause. Created 2026-08-25; updated 2026-09-16. Current status and exact special-case scope require independent verification.
\item Registry curation by Ulam AI / UnsolvedMath Contributors, CC BY 4.0; the underlying report retains its supplied provenance. Unreviewed statement-only draft.
\end{itemize}
\end{document}
